\noindent
Let $A$ be a ring and let $f_1, \ldots, f_r$ be a sequence of elements
of $A$. We have defined the Koszul complex
$K_\bullet(f_1, \ldots, f_r)$ in
More on Algebra, Definition \ref{more-algebra-definition-koszul-complex}.
It is a chain complex sitting in degrees $r, \ldots, 0$.
We turn this into a cochain complex $K^\bullet(f_1, \ldots, f_r)$
by setting $K^{-n}(f_1, \ldots, f_r) = K_n(f_1, \ldots, f_r)$
and using the same differentials. In the rest of this section all
the complexes will be cochain complexes.

\medskip\noindent
We define a complex $I^\bullet(f_1, \ldots, f_r)$
such that we have a distinguished triangle
$$
I^\bullet(f_1, \ldots, f_r) \to
A \to
K^\bullet(f_1, \ldots, f_r) \to
I^\bullet(f_1, \ldots, f_r)[1]
$$
in $K(A)$.
In other words, we set
$$
I^i(f_1, \ldots, f_r) =
\left\{
\begin{matrix}
K^{i - 1}(f_1, \ldots, f_r) & \text{if } i \leq 0 \\
0 & \text{else}
\end{matrix}
\right.
$$
and we use the negative of the differential on $K^\bullet(f_1, \ldots, f_r)$.
The maps in the distinguished triangle are the obvious ones. Note that
$I^0(f_1, \ldots, f_r) = A^{\oplus r} \to A$ is given by
multiplication by $f_i$ on the $i$th factor.
Hence $I^\bullet(f_1, \ldots, f_r) \to A$ factors as
$$
I^\bullet(f_1, \ldots, f_r) \to I \to A
$$
where $I = (f_1, \ldots, f_r)$. In fact, there is a short exact sequence
$$
0 \to H^{-1}(K^\bullet(f_1, \ldots, f_s)) \to
H^0(I^\bullet(f_1, \ldots, f_s)) \to I \to 0
$$
and for every $i < 0$ we have
$H^i(I^\bullet(f_1, \ldots, f_r)) = H^{i - 1}(K^\bullet(f_1, \ldots, f_r))$.
Observe that given a second sequence $g_1, \ldots, g_r$ of elements of $A$
there are canonical maps
$$
I^\bullet(f_1g_1, \ldots, f_rg_r) \to I^\bullet(f_1, \ldots, f_r)
\quad\text{and}\quad
K^\bullet(f_1g_1, \ldots, f_rg_r) \to K^\bullet(f_1, \ldots, f_r)
$$
compatible with the maps described above. The first of these maps is
given by multiplication by $g_i$ on the $i$th summand of
$I^0(f_1g_1, \ldots, f_rg_r) = A^{\oplus r}$. In particular, given
$f_1, \ldots, f_r$ we obtain an inverse system of complexes
\begin{equation}
\label{equation-system}
I^\bullet(f_1, \ldots, f_r) \leftarrow
I^\bullet(f_1^2, \ldots, f_r^2) \leftarrow
I^\bullet(f_1^3, \ldots, f_r^3) \leftarrow \ldots
\end{equation}
which will play an important role in that which is to follow.
To easily formulate the following lemmas we fix some notation.
